\chapter{Modelagem de Casos de Uso e Análise Conceitual}
\label{chap:modelagem_casos_uso}

\section{Identificação e Estratégia dos Atores}
Atores representam os papéis desempenhados por usuários humanos ou sistemas externos que interagem diretamente com as fronteiras do sistema.

\begin{table}[htbp]
\caption{Quadro de Atores do Sistema}
\label{tab:atores_sistema}
\centering
\small
\begin{tabularx}{\textwidth}{@{} l l Y @{}}
\toprule
\textbf{Ator} & \textbf{Tipo} & \textbf{Descrição do Papel} \\
\midrule
Recepcionista & Humano & Realiza agendamentos, confirma atendimentos e cadastra clientes. \\
Profissional Especialista & Humano & Preenche laudos, registros técnicos e atende demandas especializadas. \\
Administrador & Humano & Gerencia permissões, audita logs e configura parâmetros do sistema. \\
Gateway de Pagamento & Sistema Externo & Processa cobranças, emite confirmações de transação financeira. \\
\bottomrule
\end{tabularx}
\end{table}

\section{Diagrama Geral de Casos de Uso}
A visão geral de casos de uso delimita as funcionalidades disponibilizadas aos diferentes atores. A modelagem detalhada é particionada por módulos funcionais para evitar diagramas monolíticos com cruzamento excessivo de linhas.

\begin{center}
    \fbox{\parbox{0.9\textwidth}{\centering\vspace{1cm}\small\textbf{[Diagrama Geral de Casos de Uso]}\\\textit{Insira o arquivo em Imagens/ e utilize o comando \texttt{\textbackslash incluirdiagrama\{Imagens/nome\_arquivo\}\{Legenda\}\{label\}}}\vspace{1cm}}}
\end{center}

\section{Especificação Textual dos Casos de Uso}

\subsection{UC01 -- Realizar Agendamento}
\begin{table}[htbp]
\caption{Especificação Textual do Caso de Uso UC01}
\label{tab:uc01_especificacao}
\centering
\small
\begin{tabularx}{\textwidth}{@{} l Y @{}}
\toprule
\multicolumn{2}{@{}l}{\textbf{UC01 -- Realizar Agendamento de Atendimento}} \\
\midrule
\textbf{Atores} & Recepcionista (iniciador), Cliente. \\
\textbf{Objetivo} & Registrar uma nova sessão de atendimento para o cliente no calendário do profissional. \\
\textbf{Pré-condições} & O cliente e o profissional devem estar previamente cadastrados e ativos no sistema. \\
\textbf{Pós-condições} & O horário é bloqueado na agenda e uma notificação de confirmação é enviada. \\
\midrule
\multicolumn{2}{@{}l}{\textbf{Fluxo Principal}} \\
\midrule
\multicolumn{2}{@{}X}{
1. A recepcionista acessa a opção de agendamento no sistema. \newline
2. O sistema solicita a identificação do cliente e a especialidade desejada. \newline
3. A recepcionista informa o CPF do cliente e seleciona o profissional. \newline
4. O sistema exibe a grade de horários disponíveis do profissional selecionado. \newline
5. A recepcionista seleciona a data e horário desejados. \newline
6. O sistema valida a inexistência de conflitos de horário (RN01). \newline
7. O sistema solicita a confirmação dos dados da reserva. \newline
8. A recepcionista confirma a operação. \newline
9. O sistema persiste o agendamento com status ``Agendado'' e exibe mensagem de sucesso.
} \\
\midrule
\multicolumn{2}{@{}l}{\textbf{Fluxos Alternativos e de Exceção}} \\
\midrule
\multicolumn{2}{@{}X}{
\textbf{FA01 -- Cliente Não Cadastrado (Passo 3):} O sistema permite redirecionar para o cadastro rápido do cliente e retomar o agendamento.\newline
\textbf{FE01 -- Horário Indisponível / Conflito Simultâneo (Passo 6):} Se outro operador reservar o horário simultaneamente, o sistema informa a colisão e reapresenta os horários livres atualizados.
} \\
\midrule
\textbf{Regras de Negócio} & RN01 (Conflito de horários), RN03 (Antecedência mínima de cancelamento). \\
\bottomrule
\end{tabularx}
\end{table}

\section{Matriz de Rastreabilidade (Requisitos vs Casos de Uso)}
\begin{table}[htbp]
\caption{Matriz de Rastreabilidade entre Requisitos Funcionais e Casos de Uso}
\label{tab:rastreabilidade_rf_uc}
\centering
\small
\begin{tabularx}{\textwidth}{@{} l Y l @{}}
\toprule
\textbf{Requisito Funcional} & \textbf{Descrição Sintética} & \textbf{Casos de Uso Associados} \\
\midrule
RF01 & Autenticação e Controle de Acesso & UC00 (Autenticar Usuário) \\
RF02 & Agendamento de Consultas & UC01 (Realizar Agendamento), UC02 (Cancelar Agendamento) \\
RF03 & Prontuário e Evolução & UC03 (Registrar Atendimento Clínico) \\
RF04 & Controle de Estoque & UC04 (Dispensar Insumos), UC05 (Repor Estoque) \\
\bottomrule
\end{tabularx}
\end{table}

\section{Diagrama de Classes de Análise / Conceitual de Domínio}
O modelo de domínio descreve os conceitos e entidades fundamentais de negócio sem detalhes técnicos de implementação (banco de dados, linguagens de programação ou frameworks).
